\chapter[\emj{🌐}~Floyd-Warshall (Todos los Pares)]{\emj{🌐}~Floyd-Warshall (Caminos más Cortos entre Todos los Pares)}

\begin{dsaquees}

Una tabla de distancias de un mapa de carreteras 🗺️, donde cada casilla
te dice cuánto cuesta ir de la ciudad X a la ciudad Y.

Floyd-Warshall llena esa tabla probando algo simple para cada par:
"¿me conviene ir directo de X a Y, o es más barato pasar por una ciudad
intermedia K?". Prueba TODAS las ciudades intermedias posibles, una por
una, y se queda con la ruta más barata.

Al terminar, tienes la distancia mínima entre cualquier par de ciudades.

\end{dsaquees}

\begin{dsaotraforma}

Floyd-Warshall calcula la distancia más corta entre TODOS los pares de
nodos de un grafo dirigido y ponderado. Admite pesos negativos (pero no
ciclos negativos).

Es programación dinámica pura: para cada nodo intermedio $k$, mejora la
distancia de todo par $(i, j)$ preguntando si pasar por $k$ es más corto:
\[ d[i][j] = \min\big(d[i][j],\ d[i][k] + d[k][j]\big) \]
donde $d$ es la matriz de distancias.

Su elegancia es un triple bucle de 3 líneas. Su precio es $O(V^3)$ en
tiempo y $O(V^2)$ en memoria, así que solo sirve para grafos pequeños o
densos (cientos de nodos, no millones).

\end{dsaotraforma}

\begin{dsadiagrama}
\begin{verbatim}
Matriz de distancias (inf = sin conexión directa):

        A    B    C
   A  [ 0    3    inf ]
   B  [ inf  0    1   ]
   C  [ 2    inf  0   ]

Probando intermedio k = B:
  A->C ahora puede pasar por B: A->B(3) + B->C(1) = 4
  dist[A][C] = min(inf, 4) = 4

Probando intermedio k = C:
  B->A puede pasar por C: B->C(1) + C->A(2) = 3
  dist[B][A] = min(inf, 3) = 3

Matriz final (todas las distancias mínimas):
        A    B    C
   A  [ 0    3    4 ]
   B  [ 3    0    1 ]
   C  [ 2    5    0 ]
\end{verbatim}
\end{dsadiagrama}

\begin{lstlisting}[language=C++]
CÓMO FUNCIONA
1. dist[i][j] = peso de la arista i->j (0 si i==j, inf si no hay arista).
2. Por cada nodo intermedio k (bucle externo):
     por cada par (i, j):
       dist[i][j] = min(dist[i][j], dist[i][k] + dist[k][j])
3. ¡El orden importa! k DEBE ser el bucle más externo.
4. Si algún dist[i][i] queda negativo -> hay ciclo negativo.

📝 PSEUDOCÓDIGO
──────────────────────────────────────────────────────────────

función floydWarshall(dist):   // dist es matriz V x V
    PARA k desde 0 hasta V-1:
        PARA i desde 0 hasta V-1:
            PARA j desde 0 hasta V-1:
                SI dist[i][k] + dist[k][j] < dist[i][j]:
                    dist[i][j] = dist[i][k] + dist[k][j]
    return dist

💻 CÓDIGO C++
──────────────────────────────────────────────────────────────

#include <iostream>
#include <vector>
using namespace std;

const long long INF = 1e18;

void floydWarshall(vector<vector<long long>>& dist) {
    int V = dist.size();
    for (int k = 0; k < V; k++)
        for (int i = 0; i < V; i++)
            for (int j = 0; j < V; j++)
                if (dist[i][k] < INF && dist[k][j] < INF &&
                    dist[i][k] + dist[k][j] < dist[i][j])
                    dist[i][j] = dist[i][k] + dist[k][j];
}

int main() {
    int V = 3;
    vector<vector<long long>> dist(V, vector<long long>(V, INF));
    for (int i = 0; i < V; i++) dist[i][i] = 0;
    dist[0][1] = 3; dist[1][2] = 1; dist[2][0] = 2; // aristas
    floydWarshall(dist);
    // dist[i][j] ahora tiene la distancia mínima entre cada par
    return 0;
}

⏱️ COMPLEJIDAD (Big-O)
──────────────────────────────────────────────────────────────

  Recurso     │ Costo
  ────────────┼─────────────
  Tiempo      │ O(V^3)
  Espacio     │ O(V^2)

* Ideal para V pequeño (≤ 400-500). Para V grande usa Dijkstra por nodo.
\end{lstlisting}

\begin{dsaentrevistas}

✅ ¿Cuándo usarlo? Cuando pidan distancias entre TODOS los pares y el
   grafo sea pequeño/denso. Si solo necesitas de UNA fuente, usa
   Dijkstra o Bellman-Ford.

💡 Cierre transitivo: cambia (min, +) por (OR, AND) y Floyd-Warshall te
   dice qué nodos son alcanzables desde cuáles. Mismo esqueleto.

⚠️ El orden de los bucles NO se negocia: k SIEMPRE va afuera. Si lo
   pones adentro, el resultado es incorrecto.

🔑 Ciclo negativo: revisa la diagonal. Si algún dist[i][i] < 0, existe
   un ciclo de peso negativo pasando por i.

\end{dsaentrevistas}

\begin{dsaresumen}

• Distancias más cortas entre TODOS los pares de nodos.
• Programación dinámica: intermedio k mejora cada par (i, j).
• Triple bucle con k por fuera; $O(V^3)$ tiempo, $O(V^2)$ espacio.
• Admite pesos negativos, detecta ciclos negativos en la diagonal.
• Solo para grafos pequeños o densos.
• Alternativa: Dijkstra por cada nodo si el grafo es disperso.
════════════════════════════════════════════════════════════════

\end{dsaresumen}
